\section{Background and Related Work}

\subsection{EPSS and Exploit-Likelihood-Based Prioritization}

The Exploit Prediction Scoring System (EPSS)~\cite{jacobs2021epss} is a machine-learning model that estimates the probability that a published CVE will be exploited in the wild within the next 30 days. EPSS v3 uses $\sim$1,500 features drawn from CVE descriptors, NVD metadata, and threat intelligence feeds, and produces daily probability scores for all published CVEs. Empirical evaluations~\cite{jacobs2020improving} have shown EPSS substantially outperforms CVSS base scores as a predictor of near-term exploitation: CVEs in the top decile of EPSS account for a disproportionate fraction of observed exploitation events. EPSS is now integrated into CISA's vulnerability prioritization guidance and is used by numerous RBVM vendors.

CISA's Known Exploited Vulnerabilities (KEV) catalog~\cite{cisa2026kev} is a complementary signal: it lists CVEs confirmed to have been exploited in the wild, with mandated remediation deadlines for federal agencies under BOD 22-01~\cite{cisa2021bod2201}. KEV membership provides strong prioritization signal but covers only a small fraction of the CVE backlog ($\sim$15,000 of the 200,000+ published CVEs as of 2026). Prior work has studied optimal combined use of EPSS and KEV~\cite{paper1vulnprio}, generally finding that KEV provides a high-precision but low-recall prioritization set, while EPSS provides broader coverage at reduced precision.

\textbf{The asset-agnostic limitation.} A structural limitation of both EPSS and KEV is that they are CVE-level signals: they score a vulnerability, not a (vulnerability, host) pair. For remediation scheduling, the unit of action is the pair: patching CVE-X on Host-A is a distinct action from patching CVE-X on Host-B, with different operational cost, risk reduction, and scheduling feasibility. EPSS provides no signal about which host carrying a given CVE should be remediated first.

\subsection{CVE-Level Host-Context Augmentation}

To our knowledge, no published academic work has systematically incorporated dynamic host-level hygiene posture (patch debt, identity exposure, telemetry freshness) into EPSS-based (host, CVE) pair prioritization. Asset criticality models weight CVE severity by static host-role classifications, producing host-adjusted scores. Exposure models incorporate network reachability, service exposure, and lateral movement risk. Commercial RBVM tools (Tenable Lumin, Qualys TruRisk, Rapid7 InsightVM) combine such signals proprietarily, but their scoring methods are not published and results are not independently reproducible.

\textbf{VulnPrio~\cite{paper1vulnprio}} is the closest public predecessor to HygienePrio. It built an open, reproducible benchmark for (host, CVE) pair prioritization on a synthetic government fleet (EEHDA), evaluating 13 strategies combining EPSS, CVSS, KEV, asset criticality, endpoint exposure, and multi-factor composite scores across 30 seeds. The central finding was null: the proposed multi-factor model was statistically indistinguishable from EPSS-only on the primary metric (expected exploited host-days, EHD), with all inter-strategy differences falling within seed-to-seed variation. VulnPrio attributed this to uncalibrated weights on a synthetic fixture and framed the result as a foundation for future calibrated studies.

\textbf{Differentiating HygienePrio from VulnPrio.} Whereas VulnPrio found the multi-factor model statistically indistinguishable from EPSS-only, we hypothesize that the missing signal is host-level hygiene posture --- not additional CVE-level features alone. VulnPrio's augmentation features (asset criticality, endpoint exposure, KEV flags) are all CVE-level or static host-role features. HygienePrio adds dynamic, telemetry-derived hygiene signals: how much patch debt is currently outstanding on this specific host, how exposed this host is through AD group memberships and privileged accounts, and how recently this host has reported telemetry. These are fundamentally different in kind from asset criticality scores and network exposure flags.

\subsection{Cyber-Hygiene Benchmarks and Anomaly Detection}

\textbf{HygieneBench~\cite{paper3hygienebench}} introduced the first open benchmark for cyber-hygiene anomaly detection covering AD identity state, endpoint patch posture, vulnerability exposure, and telemetry freshness. Its evaluation of eight detection methods across 12 anomaly classes and seven tasks found that 86.2\% of (condition, task, method) configurations fail to outperform a rule baseline, but identified task T5 (patch-vulnerability hygiene) and T2 (group-membership drift) as exception zones where structured methods add meaningful signal. HygienePrio draws on the HygieneBench schema as the source of HRS features and inherits its fleet generation logic.

\textbf{Anomaly detection for hygiene monitoring} has received limited attention in the academic literature. Outlier ensemble methods~\cite{aggarwal2015outlier} address anomaly detection in tabular data but are not designed for the temporal, multi-table, identity-centric character of enterprise hygiene state. To our knowledge, no existing public benchmark simultaneously models AD identity state, endpoint patch posture, vulnerability exposure, and telemetry freshness as first-class evaluation axes; adjacent benchmarks are surveyed in~\cite{paper3hygienebench}.

\subsection{Patch Management and Remediation Scheduling}

NIST SP 800-40 Rev.~4~\cite{nist2022sp80040r4} frames enterprise patch management as a risk-based discipline requiring current, accurate asset and patch telemetry. CISA BOD 23-01~\cite{cisa2022bod2301} mandates 14-day asset discovery and 72-hour patch data freshness cadences for federal agencies. These mandates implicitly encode hygiene requirements: fleets that fail them have stale telemetry, reducing the reliability of any signal-based prioritization. HygienePrio models this explicitly through the TelemetryFreshness dimension of HRS.

Remediation scheduling under capacity constraints has been studied as an optimization problem~\cite{paper1vulnprio}. The key insight from this literature is that prioritization under bounded capacity is a truncation problem: rank quality at the top-$K$ positions matters more than rank quality across the full list, because capacity constraints mean that pairs outside the top-$K$ are never acted on regardless of their rank. This motivates Precision@$K$ as the primary evaluation metric.
